% PAKDD 2027 Applied Data Science Track submission.
% Compile: llncs.cls from Springer's LNCS author kit (llncs.zip), pdflatex.
% Double-blind: keep author block anonymous until acceptance.
% All numbers come from artifacts/knowledge-retrieval-eval/eval_306_snapshot.json,
% bm25_ablation.json, significance.json, and ontology-evidence-audit/latest.json.
\documentclass[runningheads]{llncs}

\usepackage[T1]{fontenc}
\usepackage{graphicx}
\usepackage{booktabs}
\usepackage{multirow}
\usepackage{amsmath}
\usepackage{url}
\urldef{\mailsa}\path|author-one@example.org|   % TODO: anonymous contact
\newcommand{\best}[1]{\textbf{#1}}

\title{Ask the Regulations: A Type-Stratified Study of Hybrid Retrieval
over Small Domain Corpora with Evidence-Audited Ontologies}

\titlerunning{Ask the Regulations: Type-Stratified Hybrid Retrieval with Evidence-Audited Ontologies}

\author{Anonymous Submission \#XXXX}
\authorrunning{Anonymous Submission}
\institute{Paper ID XXXX}

\begin{document}

\maketitle

\begin{abstract}
Retrieval-augmented systems deployed in regulated domains face a reality rarely
discussed in benchmark papers: the corpus is small and grows organically, queries
fall into sharply different types, and every answer must be traceable to
authoritative text. We study this regime on a corpus of 27 Chinese marine-dumping
regulatory documents and operational assets from the national environmental
authority. We first construct a domain ontology (100 nodes, 200 relations) in
which every LLM-extracted node and relation must cite its supporting passage,
and we audit the evidence chain with deterministic string verification: 96\% of
node names and 85\% of relations are fully traceable to the cited source text,
and 89.9\% of citations verify verbatim. Second, we build a 306-question
type-stratified benchmark (title-lookup, entity/content, multi-hop, refusal)
with deterministically constructed gold answers. Third, we compare six
retrieval configurations with per-question paired Wilcoxon tests and find that
\emph{no single configuration wins across query types}: plain BM25 is
significantly best on title queries; graph-expanded lexical retrieval is
decisively best on multi-hop queries (47.5\% vs.\ 7.5\% hit@5, $p<10^{-5}$);
and a strong dense retriever with a cross-encoder reranker---often assumed
superior---is significantly worst across the board on this corpus of
near-duplicate official documents. An oracle that routes each query type to its
best configuration bounds the achievable gain. For abstention, score thresholds
that separate answerable from unanswerable queries under one configuration
invert or overlap under others, cautioning against score-based refusal in
small-corpus deployments. We release the benchmark, audit tooling, and all
configurations.

\keywords{Domain retrieval \and Hybrid retrieval \and Knowledge ontology \and
Evidence audit \and Retrieval evaluation \and Abstention}
\end{abstract}

\section{Introduction}
\label{sec:intro}

Most retrieval research is evaluated on large, heterogeneous corpora where
dense retrieval and cross-encoder reranking are known to excel~\cite{beir,dpr}.
The deployment reality in regulated domains---environmental compliance,
permits, technical standards---is almost the opposite. The corpus is small
(tens of documents, not millions) and grows one announcement at a time;
documents are near-duplicates of each other, sharing boilerplate and titles
that differ by a single edition number; queries fall into sharply different
types, from title lookup to multi-hop reasoning across related instruments;
and, because a wrong citation can mislead a real governance decision, every
answer must be traceable to authoritative text. Small-corpus retrieval is not
merely large-corpus retrieval with fewer documents: score distributions,
ranking granularity, and the behaviour of fusion strategies all change.

This paper reports a controlled, type-stratified study of retrieval in exactly
this regime, on a real deployment: a knowledge service that backs a marine
debris governance agent for aquaculture zones. The service ingests public
regulatory documents from the national environmental authority (marine dumping
area announcements, standards, permit rules, sewage outfall guidelines, and
comprehensive governance plans), maintains a domain ontology, and answers
queries with citation-backed evidence.

We address three research questions:

\begin{description}
\item[RQ1] On a small, near-duplicate regulatory corpus, how do lexical,
graph-expanded, dense, and hybrid retrieval configurations compare
\emph{per query type}, with statistical significance?
\item[RQ2] Can an LLM-extracted domain ontology be made evidence-traceable by
construction, and audited deterministically without a gold standard?
\item[RQ3] Is score-threshold abstention (refusing out-of-domain queries)
viable, and is any threshold portable across configurations?
\end{description}

\paragraph{Contributions.}
\begin{enumerate}
\item \textbf{An evidence-audit framework for LLM-extracted ontologies}
(Sect.~\ref{sec:ontology}): every node and relation must cite a passage of a
registered asset; deterministic verbatim verification reports citation
verifiability (89.9\%, 302/336), relation-endpoint traceability (90.0\%), full
relation traceability (85.0\%), and node-name traceability (96\%). To our
knowledge this audit protocol---audit by exact string verification against the
cited source, no model in the loop---is itself a reusable contribution.
\item \textbf{A type-stratified 306-question benchmark} (Sect.~\ref{sec:bench})
with four query types (title, entity/content, multi-hop, refusal) and
deterministic gold construction: zero model calls, fixed templates, gold taken
from the ontology and asset tables themselves. We also document difficulty
drift as the corpus grew from 15 to 27 assets.
\item \textbf{A controlled comparison of six configurations with paired
significance tests} (Sect.~\ref{sec:results}): the headline finding is that
\emph{no configuration wins across types}; several common assumptions
(dense-plus-reranker dominance, RRF fusion as a free lunch) do not survive.
\item \textbf{A routing upper bound and an abstention fragility result}
(Sect.~\ref{sec:route}): an oracle router over query types bounds achievable
gains, and abstention thresholds exhibit three distinct behaviours across
configurations (separable, inverted, overlapping), arguing for
evidence-coverage-based refusal.
\item \textbf{Released artifacts}: benchmark, all six configurations, the
audit and significance tooling.\footnote{Anonymized repository:
\url{https://anonymous.4open.science/r/ask-the-regulations} (placeholder to be
created before submission).}
\end{enumerate}

\section{Related Work}
\label{sec:related}

\paragraph{Hybrid retrieval.} Lexical (BM25~\cite{bm25}) and dense
retrieval~\cite{dpr,realm} are commonly fused; reciprocal rank fusion
(RRF)~\cite{rrf} is a robust parameter-light combiner. BEIR~\cite{beir} showed
that no single retriever dominates across heterogeneous corpora; we show the
same holds \emph{within} one small domain across \emph{query types}, a
granularity usually hidden by corpus-level averages.

\paragraph{KG-augmented retrieval.} GraphRAG~\cite{graphrag} and
HippoRAG~\cite{hipporag} augment retrieval with graph structures extracted
from text, mostly over large corpora. We study the complementary regime: a
small, curated, versioned ontology where graph expansion is a precision
instrument for multi-hop queries---and show it \emph{hurts} title queries
relative to plain BM25.

\paragraph{LLM knowledge construction and verifiability.} LLMs are now used to
construct knowledge graphs~\cite{zhu2024llmkg,pan2024roadmap}, but extraction
hallucination is a documented failure mode~\cite{hallucination-survey}. Our
audit is a deterministic, model-free check that every extracted fact carries a
verbatim-verifiable citation---related in spirit to grounding checks for
RAG~\cite{rag-survey} but applied to ontology construction, where the unit of
verification is a triple plus its cited evidence.

\paragraph{Abstention in retrieval and RAG.} Calibration and knowing-when-not-
to-answer are studied for LLMs~\cite{lmknow,kadavath} and for retrieval-augmented
pipelines such as Self-RAG~\cite{selfrag}. We contribute a concrete negative
result: on a small regulatory corpus, top-1 retrieval scores do not yield a
portable abstention threshold, and their failure mode differs per
configuration (separation, inversion, or overlap).

\section{Corpus and Evidence-Audited Ontology}
\label{sec:ontology}

\subsection{Corpus}
The evaluation corpus mirrors a production knowledge service for marine
debris governance. It contains \textbf{27 assets}: 20 public regulatory
documents harvested from two official columns of the national environmental
authority (marine dumping-area announcements, national standards and their
draft-for-comment notices, permit rules, sewage outfall technical guidelines,
and a regional comprehensive governance plan), plus 7 operational demo assets
(event records, telemetry, monitoring datasets) that the service also indexes.
Each asset is registered once and versioned immutably.

\subsection{Ontology construction with mandatory citations}
The ontology (100 nodes, 200 relations over the current published version) is
built through a candidate-extraction, human-review, publish lifecycle. The
extraction layer mixes rule-based extraction (167 relations) and LLM
extraction (33 relations). The invariant that makes the ontology auditable is
\emph{mandatory evidence}: every node and every relation must carry at least
one evidence item $\langle$asset\_id, asset\_version\_id, citation$\rangle$,
where \emph{citation} is a snippet of the cited asset's text. Ontology
publishing records the version id, so downstream decisions can pin an ontology
version.

\subsection{Deterministic evidence audit}
\label{sec:audit-method}
The audit verifies, per relation: (i) each cited snippet appears \emph{verbatim}
(normalized whitespace) in the cited asset's registered text; and (ii) both
endpoint names occur in the text of a successfully verified citing document.
Per node: the node name (or an alias) occurs in its source document. All checks
are exact string operations over the same registered corpus---no model calls,
no manual annotation claims. We deliberately call this an
\emph{evidence-traceability audit}, not a gold-standard precision/recall
evaluation: a fully traceable relation can still be domain-silly, and recall
of ``all true facts'' is undefined without expert annotation.

\begin{table}[t]
\caption{Evidence audit of the published ontology (27 assets, 100 nodes,
200 relations).}
\label{tab:audit}
\centering
\small
\begin{tabular}{lrr}
\toprule
Audit metric & Result & Rate \\
\midrule
Citations verbatim-verifiable & 302/336 & 89.9\% \\
Relation endpoints traceable & 180/200 & 90.0\% \\
Relations fully traceable (citations + endpoints) & 170/200 & 85.0\% \\
Node names traceable to source doc & 96/100 & 96.0\% \\
\midrule
Extraction method & 167 rule / 33 LLM & --- \\
\bottomrule
\end{tabular}
\end{table}

Table~\ref{tab:audit} shows the audit outcome. The 11\% non-verbatim
citations are concentrated in LLM-extracted relations whose snippets
paraphrase rather than quote; they are flagged, not silently accepted. The
untraceable sample is logged with reasons (\emph{citation-not-in-doc},
\emph{endpoint-missing-in-citing-doc}) and feeds the human review queue.

\section{A Type-Stratified Benchmark}
\label{sec:bench}

\begin{table}[t]
\caption{Benchmark composition. Gold answers are constructed
deterministically from the ontology and asset tables; zero model calls.}
\label{tab:bench}
\centering
\small
\begin{tabular}{lrrl}
\toprule
Type & $n$ & Gold & Construction \\
\midrule
Title lookup (single) & 81 & 1 asset & 3 fixed templates over asset titles \\
Entity/content (single2) & 100 & 1 asset & 2 fixed templates per reviewed \\
 & & & ontology entity; no title words \\
Multi-hop & 120 & source assets & 2-hop chains $A{\to}B{\to}C$ from \\
 & & of chain & published relations \\
Refusal & 5 & none & out-of-domain questions \\
\bottomrule
\end{tabular}
\end{table}

The benchmark has 306 questions in four types (Table~\ref{tab:bench}),
motivated by the traffic a governance knowledge service actually sees:
\emph{title lookup} (``what does the notice on X require?''), where the title
is the strongest lexical signal; \emph{entity/content queries}, which name an
entity or requirement but never use title words, forcing content-level
matching; \emph{multi-hop queries} derived from two-hop relation chains, whose
gold assets are the sources of the chain entities---no single document
mentions both endpoints; and \emph{refusal queries} (out-of-domain), for which
the correct behaviour is to abstain.

Two design choices matter for reproducibility. First, templates are fixed and
gold answers are read from published ontology/asset tables, so the benchmark
is regenerable deterministically. Second, the benchmark was built twice as the
corpus grew (15 assets $\to$ 27 assets); we report the 27-asset version
throughout and note that absolute numbers are not comparable across corpus
sizes---difficulty rises sharply when 13 near-duplicate dumping/outfall
announcements join the pool.

\section{Retrieval Configurations}
\label{sec:configs}

All configurations share the corpus representation: assets are chunked at 800
characters, each chunk prefixed with the asset title (capped at 1600
characters). Table~\ref{tab:configs} summarizes the six systems.

\begin{table}[t]
\caption{The six configurations.}
\label{tab:configs}
\centering
\small
\begin{tabular}{lp{8.2cm}}
\toprule
Name & Description \\
\midrule
BM25 & BM25 over chunks ($k_1{=}1.5$, $b{=}0.75$); CJK text tokenized as
character bigrams, ASCII runs as whole tokens; no graph expansion. \\
Lex+Graph & Production retriever: weighted lexical scoring (title $\times$6,
body $\times$1.5, standard-number $\times$2) plus one-hop ontology graph
expansion. \\
Dense+Rerank & BGE-M3 embeddings; cosine top-30 chunks; BGE-Reranker-V2-M3
reranks top-20 chunks; asset rank = its best chunk. \\
RRF Hybrid & RRF ($k{=}60$) fusion of Lex+Graph asset ranking and dense
asset ranking. \\
TB Hybrid & RRF hybrid whose dense channel embeds title-boosted chunks
(title repeated $3\times$). \\
Oracle & Upper bound only: routes multi-hop to Lex+Graph, everything else to
TB Hybrid. \\
\bottomrule
\end{tabular}
\end{table}

BM25 uses our own index (deterministic tokenizer, no segmenter dependency);
Lex+Graph is the deployed production service; Dense+Rerank and the hybrids run
in the evaluation harness over identical chunks. All model calls are cached;
every number in this paper is recomputable from the frozen run artifact.

\section{Results}
\label{sec:results}

\begin{table}[t]
\caption{Main results on the 306-question benchmark (frozen run). Per type:
hit@1 / hit@5 / MRR over ranked assets. Bold = significantly best in the
column (paired Wilcoxon, $p<0.05$, see Table~\ref{tab:sig}).}
\label{tab:main}
\centering
\small
\setlength{\tabcolsep}{3.2pt}
\begin{tabular}{l ccc ccc ccc}
\toprule
& \multicolumn{3}{c}{Title ($n{=}81$)} & \multicolumn{3}{c}{Entity ($n{=}100$)} & \multicolumn{3}{c}{Multi-hop ($n{=}120$)} \\
\cmidrule(lr){2-4}\cmidrule(lr){5-7}\cmidrule(lr){8-10}
Config & h@1 & h@5 & MRR & h@1 & h@5 & MRR & h@1 & h@5 & MRR \\
\midrule
BM25 & \best{.185} & \best{.457} & \best{.302} & .140 & .180 & .164 & .075 & .075 & .135 \\
Lex+Graph & .123 & .432 & .254 & \best{.130} & .310 & \best{.229} & \best{.458} & \best{.475} & \best{.469} \\
Dense+Rerank & .074 & .222 & .138 & .010 & .190 & .091 & .075 & .458 & .267 \\
RRF Hybrid & .173 & .432 & .284 & .030 & .380 & .186 & .075 & .458 & .198 \\
TB Hybrid & .148 & .444 & .278 & .030 & \best{.410} & .186 & .075 & .458 & .198 \\
Oracle & .148 & .444 & .278 & .030 & \best{.410} & .186 & \best{.458} & \best{.475} & \best{.469} \\
\bottomrule
\end{tabular}
\end{table}

\begin{table}[t]
\caption{Selected paired Wilcoxon signed-rank tests over per-question
reciprocal rank (two-sided; zeros dropped). Full matrix in the artifact.}
\label{tab:sig}
\centering
\small
\begin{tabular}{llrl}
\toprule
Type & Pair & $p$ & Reading \\
\midrule
Multi-hop & Lex+Graph vs.\ BM25 & $3.3{\times}10^{-6}$ & graph expansion $+40$pt hit@5 \\
Multi-hop & Lex+Graph vs.\ RRF Hybrid & $2.3{\times}10^{-10}$ & RRF fusion \emph{hurts} graph ranking \\
Title & BM25 vs.\ Lex+Graph & $1.5{\times}10^{-2}$ & plain BM25 beats graph expansion \\
Title & BM25 vs.\ Dense+Rerank & $1.6{\times}10^{-4}$ & dense channel worst on titles \\
Entity & TB Hybrid vs.\ Lex+Graph & $2.99{\times}10^{-1}$ & not significant \\
Entity & BM25 vs.\ Dense+Rerank & $9.1{\times}10^{-1}$ & not significant \\
\bottomrule
\end{tabular}
\end{table}

\paragraph{F1: No configuration wins across query types.}
Table~\ref{tab:main} is the core result. On title queries, plain BM25 is
significantly best (hit@5 .457; vs.\ Lex+Graph $p{=}.015$, vs.\ Dense+Rerank
$p{=}1.6{\times}10^{-4}$): title matching is exactly what BM25's term
weighting does, and graph expansion only adds noise. On multi-hop queries the
picture inverts: Lex+Graph reaches .475 hit@5 while BM25 collapses to .075
($p{=}3.3{\times}10^{-6}$)---the gold documents share no vocabulary with the
query, and only graph traversal bridges them. On entity/content queries the
title-boosted hybrid is nominally best (.410 hit@5) but the gap to Lex+Graph
is not significant ($p{=}.30$): content matching is the least separated axis
at this corpus size.

\paragraph{F2: Dense retrieval with cross-encoder reranking can be
significantly worst.} The configuration most often recommended by default---
bi-encoder retrieval plus a cross-encoder---is significantly worst on two of
three types (title: hit@1 .074; entity: hit@1 .010). The mechanism is
structural: this corpus contains clusters of near-duplicate official
documents; whole-chunk embeddings dilute the few discriminative tokens (an
edition year, a place name), and the reranker can only reorder the top-30
candidate chunks the retriever already ranked. This is a cautionary result for
small-domain deployments: dense retrieval's known weakness on
near-duplicates~\cite{beir} dominates when the \emph{entire corpus} is
near-duplicates of a few families.

\paragraph{F3: Naive fusion is not a free lunch.} RRF fusion of Lex+Graph
with the dense ranking \emph{preserves} title-type performance but
\emph{damages} the multi-hop sweet spot: multi-hop hit@1 drops from .458
(Lex+Graph) to .075 (RRF Hybrid; $p{=}2.3{\times}10^{-10}$), because the dense
channel's noisy ranking of bridge documents leaks into the fused list.
Title-boosting the dense channel (TB Hybrid) repairs title/entity behaviour
and is the best deployable single configuration overall, but still inherits
the fused multi-hop weakness.

\paragraph{F4: Difficulty drift is real.} On the earlier 15-asset corpus the
same Lex+Graph configuration scored 1.000 hit@5 on multi-hop queries and the
dense two-stage pipeline led entity queries by $+33$ points; after 12 more
near-duplicate announcements entered the corpus, every absolute number fell
and the ordering \emph{within} types changed. Small-corpus leaderboards age
quickly; type-stratified reporting makes the drift visible instead of hiding
it in an average.

\section{Routing Upper Bound and Abstention}
\label{sec:route}

\paragraph{Oracle routing.} Routing multi-hop queries to Lex+Graph and all
others to TB Hybrid yields the best of both (Table~\ref{tab:main}, last row):
it matches Lex+Graph on multi-hop and TB Hybrid elsewhere, dominating every
single configuration. The oracle knows the query type; a deployable system
must predict it. At 306 questions we deliberately do not train a classifier
---we report the bound and note that query-type classification (or a cheap
router signal such as ``does the query match any asset title?'') is the
obvious next step; even the naive fixed prior ``route title-like queries to
BM25, the rest to Lex+Graph'' already beats any single system.

\begin{table}[t]
\caption{Abstention: top-1 retrieval scores of answerable vs.\ refusal
queries. Three behaviours across configurations.}
\label{tab:refuse}
\centering
\small
\begin{tabular}{lrrl}
\toprule
Config & min answerable & max refusal & behaviour \\
\midrule
Lex+Graph & 67.5 & 46.5 & separable \\
BM25 & 9.05 & 13.72 & \emph{inverted} (refusal scores higher) \\
Dense+Rerank & 0.0008 & 0.0031 & overlap \\
\bottomrule
\end{tabular}
\end{table}

\paragraph{Abstention is configuration-fragile.} With only 5 refusal queries
this is a case study, not a powered experiment---but the pattern is stark
(Table~\ref{tab:refuse}). Under Lex+Graph, answerable and refusal scores
separate cleanly. Under BM25 the order \emph{inverts}: out-of-domain
questions score \emph{higher} than the weakest answerable ones (BM25 rewards
long out-of-domain queries that match high-IDF bigrams). Under dense scores
the two ranges overlap. The practical consequence: a cutoff tuned on one
configuration transfers to none other, and the failure is silent. Our
production QA layer therefore abstains on \emph{evidence coverage} (are the
retrieved citations sufficient to ground an answer?) rather than retrieval
score, and states its refusal to the user.

\section{Discussion and Limitations}
\label{sec:limits}

\emph{Single domain, single language.} The corpus is Chinese regulatory text;
the benchmark templates and the near-duplicate structure are domain-shaped.
Replication on a second regulatory domain is future work; the construction
recipe (Sect.~\ref{sec:bench}) is domain-agnostic.

\emph{Refusal pool is small} (5 queries), so Table~\ref{tab:refuse} is
evidence of fragility, not a powered estimate of abstention error rates.

\emph{Run-to-run variance on entity queries.} Repeat runs with an evolving
embedding cache moved entity-type absolute numbers by up to $\pm$0.11 while
preserving relative order; we therefore anchor claims on within-run paired
tests, not cross-run decimals.

\emph{Audit is not a gold standard.} Traceability verifies that extracted
knowledge quotes its source, not that the knowledge base is complete or
domain-correct (Sect.~\ref{sec:audit-method}).

\emph{Oracle routing needs a real router.} The bound assumes query types are
known; a lightweight, robust router is the missing piece between this study
and a deployable adaptive retriever.

\section{Conclusion}
\label{sec:conclusion}

On a small, near-duplicate regulatory corpus---the regime real governance
deployments live in---no retrieval configuration wins across query types:
plain BM25, graph-expanded lexical retrieval, and title-boosted hybrid
retrieval each own one query type, and a strong dense-plus-reranker pipeline
is significantly worst. Score-based abstention does not transfer across
configurations. The practical recipe our deployment adopted is: keep the
ontology evidence-audited, keep BM25-honest title matching, route by query
type when possible, and abstain on evidence coverage rather than retrieval
score. We release everything needed to re-run this study on a new domain.

\subsubsection*{Reproducibility.} The frozen run artifact, benchmark,
ablation, audit, and significance scripts are in the anonymized repository
(footnote 1). Every model call is cached; all numbers regenerate
deterministically from the frozen snapshot.

%
% Bibliography (llncs numeric style)
%
\begin{thebibliography}{25}

\bibitem{bm25}
Robertson, S., Zaragoza, H.: The probabilistic relevance framework: BM25 and
beyond. Found.\ Trends Inf.\ Retr.\ \textbf{3}(4), 333--389 (2009)

\bibitem{dpr}
Karpukhin, V., Oguz, B., Min, S., Lewis, P., Wu, L., Edunov, S., Chen, D.,
Yih, W.: Dense passage retrieval for open-domain question answering. In:
Proc.\ EMNLP (2020)

\bibitem{realm}
Guu, K., Lee, K., Tung, Z., Pasupat, P., Chang, M.: REALM: Retrieval-augmented
language model pre-training. In: Proc.\ ICML (2020)

\bibitem{rrf}
Cormack, G.V., Clarke, C.L.A., Buettcher, S.: Reciprocal rank fusion
outperforms Condorcet and individual rank learning methods. In: Proc.\ SIGIR
(2009)

\bibitem{beir}
Thakur, N., Reimers, N., R{\"u}ckl{\'e}, A., Srivastava, A., Gurevych, I.:
BEIR: A heterogeneous benchmark for zero-shot evaluation of information
retrieval models. In: Proc.\ NeurIPS Datasets and Benchmarks (2021)

\bibitem{graphrag}
Edge, D., Trinh, H., Cheng, N., et al.: From local to global: A Graph RAG
approach to query-focused summarization (2024). arXiv:2404.16130

\bibitem{hipporag}
Guti{\'e}rrez, B.J., Shu, Y., Gu, Y., Yasunaga, M., Su, Y.: HippoRAG:
Neurobiologically inspired long-term memory for large language models. In:
Proc.\ NeurIPS (2024)

\bibitem{zhu2024llmkg}
Zhu, Y., Wang, X., Chen, J., et al.: LLMs for knowledge graph construction and
reasoning: recent capabilities and future opportunities. World Wide Web
\textbf{27}, 58 (2024)

\bibitem{pan2024roadmap}
Pan, S., Luo, L., Wang, Y., et al.: Unifying large language models and
knowledge graphs: A roadmap. IEEE TKDE \textbf{36}(7), 3580--3599 (2024)

\bibitem{hallucination-survey}
Ji, Z., Lee, N., Frieske, R., et al.: Survey of hallucination in natural
language generation. ACM Comput.\ Surv.\ \textbf{55}(12), 1--38 (2023)

\bibitem{rag-survey}
Gao, Y., Xiong, Y., Gao, X., et al.: Retrieval-augmented generation for large
language models: A survey (2023). arXiv:2312.10997

\bibitem{lmknow}
Yin, Z., Sun, J., Chang, C., Guo, Q., Dai, J., Huang, X., Qiu, X.: Do large
language models know what they do not know? In: Proc.\ Findings of ACL
(2023)

\bibitem{kadavath}
Kadavath, S., Conerly, T., Askell, A., et al.: Language models (mostly) know
what they know (2022). arXiv:2207.05221

\bibitem{selfrag}
Asai, A., Wu, Z., Wang, Y., Sil, A., Hajishirzi, H.: Self-RAG: Learning to
retrieve, generate, and critique through self-reflection. In: Proc.\ ICLR
(2024)

\bibitem{bge-m3}
Chen, J., Xiao, S., Zhang, P., Luo, Y., Lian, D., Liu, Z.: BGE M3-embedding:
Multi-lingual, multi-functionality, multi-granularity text embeddings through
self-knowledge distillation. In: Proc.\ Findings of ACL (2024)

\bibitem{lewis-rag}
Lewis, P., Perez, E., Piktus, A., et al.: Retrieval-augmented generation for
knowledge-intensive NLP tasks. In: Proc.\ NeurIPS (2020)

\end{thebibliography}

\end{document}
